\documentclass[10pt,landscape]{article}
\usepackage[utf8]{inputenc}
\usepackage[margin=0.2in]{geometry}
\usepackage{amsmath,amssymb}
\usepackage{multicol}
\usepackage{enumitem}
% Reduce whitespace and tighten layout
\setlength{\parskip}{0pt}
\setlength{\columnsep}{0.1in}
\setlength{\abovedisplayskip}{2pt}
\setlength{\belowdisplayskip}{2pt}
\setlist[itemize]{noitemsep, topsep=0pt, partopsep=0pt, parsep=0pt}
\usepackage[breakable]{tcolorbox}
\usepackage{xcolor}

% Define custom colors for each topic
\definecolor{regression}{RGB}{230,240,255}
\definecolor{classification}{RGB}{255,240,230}
\definecolor{neural}{RGB}{240,255,240}
\definecolor{svm}{RGB}{255,240,255}
\definecolor{clustering}{RGB}{245,245,220}
\definecolor{mixture}{RGB}{230,255,255}
\definecolor{dimension}{RGB}{255,230,230}
\definecolor{graphical}{RGB}{240,240,255}
\definecolor{hmm}{RGB}{255,255,230}
\definecolor{mdp}{RGB}{255,235,245}
\definecolor{rl}{RGB}{235,255,235}

\pagestyle{empty}
\setlength{\columnsep}{0.1in}
\renewcommand{\arraystretch}{1.1}

\begin{document}
\scriptsize

\begin{multicols}{3}

%%%%%%%%%%%%%%%%%%%%%%%%%%%%%%%%%%%%%%%%%%
\begin{tcolorbox}[breakable, title=1. Regression Concepts, colback=regression,boxsep=1pt,left=1pt,right=1pt,top=1pt,bottom=1pt]
\textbf{Linear Regression Overview:}
\begin{itemize}[leftmargin=*,noitemsep,topsep=0pt]
    \item Maps input $x \in \mathbb{R}^D$ to target $y \in \mathbb{R}$ using weights $w$
    \item Basic form: $y = w_0 + w_1x_1 + ... + w_Dx_D$
    \item Uses least squares loss to find optimal weights
    \item "Bias trick": Add $x_0=1$ to make notation cleaner: $y = w^T x$
\end{itemize}

\textbf{Key Considerations:}
\begin{itemize}[leftmargin=*,noitemsep,topsep=0pt]
    \item Basis functions ($\phi(x)$): Transform inputs to capture non-linear relationships
    \item Regularization: Prevent overfitting by penalizing large weights
    \item Bias-variance tradeoff: Balance between underfitting and overfitting
\end{itemize}

\textbf{Regularization Types:}
\begin{itemize}[leftmargin=*,noitemsep,topsep=0pt]
    \item Ridge (L2): $\frac{\lambda}{2}w^Tw$ - Shrinks all weights toward zero
    \item Lasso (L1): $\lambda\|w\|_1$ - Creates sparse solutions (some weights = 0)
    \item Elastic Net: Combination of L1 and L2
\end{itemize}

\textbf{Bayesian View:}
\begin{itemize}[leftmargin=*,noitemsep,topsep=0pt]
    \item Model noise: $y_n \sim \mathcal{N}(w^Tx_n, \beta^{-1})$
    \item Prior on weights: $w \sim \mathcal{N}(0, S_0^{-1})$
    \item MAP estimation equivalent to ridge regression with $\lambda = \frac{s}{\beta}$
\end{itemize}
\end{tcolorbox}

%%%%%%%%%%%%%%%%%%%%%%%%%%%%%%%%%%%%%%%%%%
\begin{tcolorbox}[breakable, title=1.1 Ridge Regression Derivation, colback=regression,boxsep=1pt,left=1pt,right=1pt,top=1pt,bottom=1pt]
\textbf{When to use:} Use ridge regression when you suspect multicollinearity in your features or need to prevent overfitting in high-dimensional data.

\textbf{Objective:} Find $w_{ridge}$ that minimizes OLS with L2 regularization.

1. Start with ridge loss function:
\begin{align}
L(w) = \frac{1}{2}\sum_{n=1}^{N}(y_n - w^T\phi_n)^2 + \frac{\lambda}{2}w^Tw
\end{align}

2. In matrix form with design matrix $\Phi$ and target $y$:
\begin{align}
L(w) = \frac{1}{2}(y - \Phi w)^T(y - \Phi w) + \frac{\lambda}{2}w^Tw
\end{align}

3. Take gradient, set to zero:
\begin{align}
\nabla L(w) &= -\Phi^Ty + \Phi^T\Phi w + \lambda w = 0\\
\Rightarrow (\Phi^T\Phi + \lambda I)w &= \Phi^Ty
\end{align}

4. Optimal weights: $w_{ridge} = (\Phi^T\Phi + \lambda I)^{-1}\Phi^Ty$

\textbf{Comparison with OLS:} 
OLS solution is $w^* = (\Phi^T\Phi)^{-1}\Phi^Ty$. 
Ridge adds $\lambda I$ to $\Phi^T\Phi$, which:
\begin{itemize}[leftmargin=*,noitemsep,topsep=0pt]
    \item Makes matrix inversion more stable
    \item Shrinks weights toward zero
    \item Reduces overfitting by introducing bias
\end{itemize}
\end{tcolorbox}

%%%%%%%%%%%%%%%%%%%%%%%%%%%%%%%%%%%%%%%%%%
\begin{tcolorbox}[breakable, title=1.2 Maximum Likelihood Estimation, colback=regression,boxsep=1pt,left=1pt,right=1pt,top=1pt,bottom=1pt]
\textbf{When to use:} Use MLE when you want to find parameters that make your observed data most probable under your chosen model.

\textbf{General MLE Procedure:}

1. Write likelihood function (probability of data given parameters):
$p(D|\theta) = \prod_{n=1}^{N} p(x_n|\theta)$

2. Take logarithm for easier optimization:
$\log p(D|\theta) = \sum_{n=1}^{N} \log p(x_n|\theta)$

3. Find parameters maximizing log-likelihood:
$\frac{\partial}{\partial \theta}\log p(D|\theta) = 0$

\textbf{Example: Linear Regression MLE}

1. Assume $y_n \sim \mathcal{N}(w^T x_n, \beta^{-1})$

2. Log-likelihood:
$\log p(y|X, w, \beta) = \frac{N}{2}\log\beta - \frac{N}{2}\log(2\pi) - \frac{\beta}{2}\sum_{n}(y_n - w^Tx_n)^2$

3. Take derivative, set to zero:
$\frac{\partial \log p}{\partial w} = \beta \sum_{n}(y_n - w^Tx_n)x_n = 0$

4. Solve with matrix notation:
$X^Ty = X^TXw \Rightarrow w^* = (X^TX)^{-1}X^Ty$

\textbf{Key insight:} MLE for linear regression with Gaussian noise is equivalent to minimizing the least squares loss!
\end{tcolorbox}

%%%%%%%%%%%%%%%%%%%%%%%%%%%%%%%%%%%%%%%%%%
\begin{tcolorbox}[breakable, title=1.3 Maximum A Posteriori Estimation, colback=regression,boxsep=1pt,left=1pt,right=1pt,top=1pt,bottom=1pt]
\textbf{When to use:} Use MAP when you have prior knowledge about parameters and want to incorporate this information into your estimation.

\textbf{MAP General Procedure:}

1. Write posterior using Bayes' rule:
$p(\theta|D) \propto p(D|\theta)p(\theta)$

2. Take logarithm:
$\log p(\theta|D) = \log p(D|\theta) + \log p(\theta) + const.$

3. Maximize log posterior:
$\frac{\partial}{\partial \theta}[\log p(D|\theta) + \log p(\theta)] = 0$

\textbf{Example: Bayesian Linear Regression}

1. Gaussian prior: $w \sim \mathcal{N}(0, S_0^{-1})$

2. Log posterior:
$\log p(w|X,y,\beta) \propto -\frac{\beta}{2}\sum_n(y_n - w^Tx_n)^2 - \frac{1}{2}w^TS_0w$

3. Take derivative, set to zero:
$\beta\sum_n(y_n - w^Tx_n)x_n - S_0w = 0$

4. Matrix form solution:
$w_{MAP} = (\beta X^TX + S_0)^{-1}\beta X^Ty$

\textbf{Note:} With $S_0 = \lambda I$, we get ridge regression: 
$w_{MAP} = (X^TX + \frac{\lambda}{\beta}I)^{-1}X^Ty$
\end{tcolorbox}

%%%%%%%%%%%%%%%%%%%%%%%%%%%%%%%%%%%%%%%%%%
\begin{tcolorbox}[breakable, title=1.4 Posterior Parameters \& Predictive Distribution, colback=regression,boxsep=1pt,left=1pt,right=1pt,top=1pt,bottom=1pt]
\textbf{When to use:} Use posterior predictive distribution when you want to make predictions that account for uncertainty in your parameters.

\textbf{Finding Posterior Parameters (Conjugate Distributions):}

1. Identify prior distribution and parameters
2. Identify likelihood function
3. Use conjugacy to determine posterior form
4. Derive posterior parameters by comparing coefficients

\textbf{Example: Gaussian-Gaussian conjugacy}
\begin{itemize}[leftmargin=*,noitemsep,topsep=0pt]
    \item Prior: $w \sim \mathcal{N}(\mu_0, \Sigma_0)$
    \item Likelihood: $y|x, w \sim \mathcal{N}(w^Tx, \sigma^2)$
    \item Posterior: $w|D \sim \mathcal{N}(\mu_N, \Sigma_N)$ where:
    \begin{align}
    \Sigma_N^{-1} &= \Sigma_0^{-1} + \frac{1}{\sigma^2}\sum_{n}x_nx_n^T\\
    \mu_N &= \Sigma_N(\Sigma_0^{-1}\mu_0 + \frac{1}{\sigma^2}\sum_{n}y_nx_n)
    \end{align}
\end{itemize}

\textbf{Posterior Predictive Distribution:}
\begin{align}
p(y^*|x^*, D) = \int p(y^*|x^*, \theta)p(\theta|D)d\theta
\end{align}

\textbf{For Gaussian posterior:}
$y^*|x^*, D \sim \mathcal{N}(\mu_N^Tx^*, \sigma^2 + x^{*T}\Sigma_Nx^*)$

\textbf{Key insight:} The posterior predictive variance includes both noise variance $\sigma^2$ and parameter uncertainty $x^{*T}\Sigma_Nx^*$.
\end{tcolorbox}

%%%%%%%%%%%%%%%%%%%%%%%%%%%%%%%%%%%%%%%%%%
\begin{tcolorbox}[breakable, title=2. Classification Concepts, colback=classification,boxsep=1pt,left=1pt,right=1pt,top=1pt,bottom=1pt]
\textbf{Classification Overview:}
\begin{itemize}[leftmargin=*,noitemsep,topsep=0pt]
    \item Maps input $x$ to discrete class labels $y$
    \item Can be binary (two classes) or multi-class
    \item Decision boundaries separate different classes
\end{itemize}

\textbf{Common Classification Methods:}
\begin{itemize}[leftmargin=*,noitemsep,topsep=0pt]
    \item Discriminant functions: Direct mapping from inputs to class decisions
    \item Probabilistic discriminative models: Model $p(y|x)$ directly (logistic regression)
    \item Probabilistic generative models: Model $p(x,y)$ together (Naive Bayes)
\end{itemize}

\textbf{Loss Functions:}
\begin{itemize}[leftmargin=*,noitemsep,topsep=0pt]
    \item 0/1 Loss: Simple but non-differentiable, counts misclassifications
    \item Logistic Loss: Smooth approximation of 0/1 loss, used in logistic regression
    \item Hinge Loss: Used in SVMs, penalizes points close to or on wrong side of margin
\end{itemize}

\textbf{Training Techniques:}
\begin{itemize}[leftmargin=*,noitemsep,topsep=0pt]
    \item Gradient Descent: Iterative optimization to minimize loss
    \item Batch GD: Updates using all training examples
    \item SGD: Updates using one example at a time, faster but noisier
\end{itemize}

\textbf{Generative vs. Discriminative:}
\begin{itemize}[leftmargin=*,noitemsep,topsep=0pt]
    \item Discriminative: Models decision boundary directly
    \item Generative: Models how data is generated, can generate new samples
    \item Naive Bayes: Assumes feature independence given class
\end{itemize}
\end{tcolorbox}

%%%%%%%%%%%%%%%%%%%%%%%%%%%%%%%%%%%%%%%%%%
\begin{tcolorbox}[breakable, title=2.1 Logistic Regression Gradients, colback=classification,boxsep=1pt,left=1pt,right=1pt,top=1pt,bottom=1pt]
\textbf{When to use:} Use logistic regression when you need probabilistic class membership for binary or multi-class classification with a linear decision boundary.

\textbf{Binary Logistic Regression:}

1. Model: $p(y=1|x) = \sigma(w^Tx) = \frac{1}{1+e^{-w^Tx}}$

2. Likelihood: $p(\{y_i\}_{i=1}^{N}|w) = \prod_{i=1}^{N}\hat{y}_i^{y_i}(1-\hat{y}_i)^{1-y_i}$
where $\hat{y}_i = \sigma(w^Tx_i)$

3. Negative log-likelihood (cross-entropy):
$E(w) = -\sum_{i=1}^{N}\{y_i\log\hat{y}_i + (1-y_i)\log(1-\hat{y}_i)\}$

4. Gradient: $\nabla E(w) = \sum_{i=1}^{N}(\hat{y}_i - y_i)x_i$

\textbf{Multi-class with Softmax:}
$\nabla_{w_j}E(W) = \sum_{i=1}^{N}(\hat{y}_{ij} - y_{ij})x_i$
where $\hat{y}_{ij} = \text{softmax}_j(Wx_i)$

\textbf{Softmax function:} $\text{softmax}_k(z) = \frac{\exp(z_k)}{\sum_{j=1}^K\exp(z_j)}$
\end{tcolorbox}

%%%%%%%%%%%%%%%%%%%%%%%%%%%%%%%%%%%%%%%%%%
\begin{tcolorbox}[breakable, title=2.2 Generative Models for Classification, colback=classification,boxsep=1pt,left=1pt,right=1pt,top=1pt,bottom=1pt]
\textbf{When to use:} Use generative models when you want to model the data distribution, generate new samples, or when you have prior knowledge about how the data is generated.

\textbf{Framework:}
\begin{enumerate}[leftmargin=*,noitemsep,topsep=0pt]
    \item Define class prior $p(C_k)$ and class-conditional $p(x|C_k)$
    \item Joint distribution: $p(x, C_k) = p(C_k)p(x|C_k)$
    \item Full log-likelihood for binary classification:
    $\log p(\pi, \mu_1, \mu_2, \Sigma) = \sum_{i=1}^{N}y_i\log[\pi\mathcal{N}(x_i|\mu_1, \Sigma)] + (1-y_i)\log[(1-\pi)\mathcal{N}(x_i|\mu_2, \Sigma)]$
    \item Maximize to find optimal parameters
\end{enumerate}

\textbf{MLE Solutions (Gaussian with shared covariance):}
\begin{align}
\pi &= \frac{N_1}{N}\\
\mu_k &= \frac{1}{N_k}\sum_{i:y_i=k}x_i\\
\Sigma &= \frac{1}{N}\sum_{i=1}^{N}[y_i(x_i - \mu_1)(x_i - \mu_1)^T \\
&\quad + (1-y_i)(x_i - \mu_2)(x_i - \mu_2)^T]
\end{align}

\textbf{Using Lagrangian Method:}
1. Identify constraints (e.g., $\sum_{k}\pi_k = 1$)
2. Form Lagrangian: 
   $\mathcal{L}(\theta, \lambda) = \log p(D|\theta) + \lambda(c - \sum_{k}\pi_k)$
3. Take derivatives w.r.t. parameters and Lagrange multipliers
4. Solve the resulting system
\end{tcolorbox}

%%%%%%%%%%%%%%%%%%%%%%%%%%%%%%%%%%%%%%%%%%
\begin{tcolorbox}[breakable, title=2.3 Naive Bayes, colback=classification,boxsep=1pt,left=1pt,right=1pt,top=1pt,bottom=1pt]
\textbf{When to use:} Use Naive Bayes when features can reasonably be assumed to be conditionally independent given the class, or when you have limited training data.

\textbf{Key Idea:} Assume features are conditionally independent given the class:
\begin{align}
p(x|y=C_k) = \prod_{d=1}^D p(x_d|y=C_k)
\end{align}

\textbf{Classification Rule:}
\begin{align}
\hat{y} = \arg\max_k p(y=C_k|x) = \arg\max_k p(y=C_k)\prod_{d=1}^D p(x_d|y=C_k)
\end{align}

\textbf{Parameter Estimation:}
\begin{itemize}[leftmargin=*,noitemsep,topsep=0pt]
    \item Class priors: $p(y=C_k) = \frac{N_k}{N}$
    \item For discrete features: $p(x_d=v|y=C_k) = \frac{\text{count}(x_d=v, y=C_k)}{\text{count}(y=C_k)}$
    \item For continuous features: Often assume Gaussian $p(x_d|y=C_k) = \mathcal{N}(\mu_{dk}, \sigma_{dk}^2)$
\end{itemize}

\textbf{Advantages:}
\begin{itemize}[leftmargin=*,noitemsep,topsep=0pt]
    \item Fast training and prediction
    \item Works well with small datasets
    \item Handles missing values naturally
\end{itemize}

\textbf{Limitations:}
\begin{itemize}[leftmargin=*,noitemsep,topsep=0pt]
    \item Independence assumption often violated
    \item Cannot learn feature interactions
\end{itemize}
\end{tcolorbox}

%%%%%%%%%%%%%%%%%%%%%%%%%%%%%%%%%%%%%%%%%%
\begin{tcolorbox}[breakable, title=2.4 Gradient Descent Implementation, colback=classification,boxsep=1pt,left=1pt,right=1pt,top=1pt,bottom=1pt]
\textbf{When to use:} Use gradient descent when direct analytical solutions aren't possible or practical, especially with non-convex or high-dimensional loss functions.

\textbf{Basic Gradient Descent:}
\begin{enumerate}[leftmargin=*,noitemsep,topsep=0pt]
    \item Initialize parameters $\theta^{(0)}$
    \item Choose learning rate $\eta$
    \item For iterations $t=0,1,2...$:
    \begin{itemize}[leftmargin=*,noitemsep]
        \item Compute gradient: $\nabla L(\theta^{(t)})$
        \item Update: $\theta^{(t+1)} = \theta^{(t)} - \eta \nabla L(\theta^{(t)})$
        \item Check convergence
    \end{itemize}
\end{enumerate}

\textbf{For Logistic Regression:}
$w^{(t+1)} = w^{(t)} - \eta \sum_{i=1}^{N}(\sigma(w^{(t)T}x_i) - y_i)x_i$

\textbf{Stochastic Gradient Descent (SGD):}
Update using one or a batch of samples:
$w^{(t+1)} = w^{(t)} - \eta \sum_{i \in \text{batch}}(\sigma(w^{(t)T}x_i) - y_i)x_i$

\textbf{Learning Rate Selection:}
\begin{itemize}[leftmargin=*,noitemsep,topsep=0pt]
    \item Too large: May diverge or oscillate
    \item Too small: Slow convergence
    \item Adaptive rates: Decrease over time or use algorithms like Adam, RMSprop
\end{itemize}

\textbf{Convergence Criteria:}
\begin{itemize}[leftmargin=*,noitemsep,topsep=0pt]
    \item Small parameter change: $\|\theta^{(t+1)} - \theta^{(t)}\| < \epsilon$
    \item Small gradient: $\|\nabla L(\theta^{(t)})\| < \epsilon$
    \item Maximum iterations reached
\end{itemize}
\end{tcolorbox}

%%%%%%%%%%%%%%%%%%%%%%%%%%%%%%%%%%%%%%%%%%
\begin{tcolorbox}[breakable, title=4. Support Vector Machines Concepts, colback=svm,boxsep=1pt,left=1pt,right=1pt,top=1pt,bottom=1pt]
\textbf{SVM Overview:}
\begin{itemize}[leftmargin=*,noitemsep,topsep=0pt]
    \item Maximum margin classifier
    \item Finds hyperplane that maximally separates classes
    \item Uses support vectors (points closest to decision boundary)
    \item Can handle linearly and non-linearly separable data
\end{itemize}

\textbf{Key Components:}
\begin{itemize}[leftmargin=*,noitemsep,topsep=0pt]
    \item Decision function: $f(x) = w^Tx + w_0$
    \item Margin: Distance from decision boundary to nearest data points
    \item Support vectors: Data points on margin boundaries
    \item Regularization parameter C: Controls trade-off between margin size and misclassification
\end{itemize}

\textbf{Types of SVMs:}
\begin{itemize}[leftmargin=*,noitemsep,topsep=0pt]
    \item Hard-margin: No misclassifications allowed, requires linearly separable data
    \item Soft-margin: Allows some misclassifications, more robust to outliers
    \item Kernel SVMs: Transform data to higher dimensions for non-linear boundaries
\end{itemize}

\textbf{Common Kernels:}
\begin{itemize}[leftmargin=*,noitemsep,topsep=0pt]
    \item Linear: $K(x,z) = x^Tz$
    \item Polynomial: $K(x,z) = (1 + x^Tz)^q$
    \item RBF (Gaussian): $K(x,z) = \exp(-\gamma\|x-z\|^2)$
\end{itemize}

\textbf{Comparison to Logistic Regression:}
\begin{itemize}[leftmargin=*,noitemsep,topsep=0pt]
    \item SVM: Maximizes margin, focuses on boundary points
    \item Logistic regression: Probabilistic outputs, influenced by all points
\end{itemize}
\end{tcolorbox}

%%%%%%%%%%%%%%%%%%%%%%%%%%%%%%%%%%%%%%%%%%
\begin{tcolorbox}[breakable, title=4.1 Maximum Margin Classification, colback=svm,boxsep=1pt,left=1pt,right=1pt,top=1pt,bottom=1pt]
\textbf{When to use:} Use hard-margin SVM when your data is linearly separable and you want a decision boundary with maximum distance to all classes.

\textbf{Hard Margin Formulation:}
\begin{align}
\min_{w,w_0} \frac{1}{2}\|w\|^2 \\
\text{subject to}~y_i(w^Tx_i + w_0) \geq 1,~\forall i
\end{align}

\textbf{Interpretation:}
\begin{itemize}[leftmargin=*,noitemsep,topsep=0pt]
    \item Decision boundary: $w^Tx + w_0 = 0$
    \item Margin width = $\frac{2}{\|w\|}$
    \item Support vectors exactly satisfy $y_i(w^Tx_i + w_0) = 1$
\end{itemize}

\textbf{Soft Margin Formulation:}
\begin{align}
\min_{w,w_0,\xi} \frac{1}{2}\|w\|^2 + C\sum_{i=1}^N \xi_i \\
\text{subject to}~y_i(w^Tx_i + w_0) \geq 1 - \xi_i,~\xi_i \geq 0,~\forall i
\end{align}

\textbf{Slack Variables $\xi_i$:}
\begin{itemize}[leftmargin=*,noitemsep,topsep=0pt]
    \item $\xi_i = 0$: Point correctly classified beyond margin
    \item $0 < \xi_i \leq 1$: Point correctly classified but inside margin
    \item $\xi_i > 1$: Point misclassified
\end{itemize}

\textbf{Parameter C:}
\begin{itemize}[leftmargin=*,noitemsep,topsep=0pt]
    \item Large C: Less regularization, small margin, fewer misclassifications
    \item Small C: More regularization, large margin, more misclassifications
\end{itemize}
\end{tcolorbox}

%%%%%%%%%%%%%%%%%%%%%%%%%%%%%%%%%%%%%%%%%%
\begin{tcolorbox}[breakable, title=4.2 The Kernel Trick, colback=svm,boxsep=1pt,left=1pt,right=1pt,top=1pt,bottom=1pt]
\textbf{When to use:} Use kernel SVMs when your data is not linearly separable in the original feature space but may be separable in a higher-dimensional space.

\textbf{Key Insight:} The SVM dual formulation only requires inner products between data points, which can be replaced with kernel functions.

\textbf{Dual Formulation:}
\begin{align}
\max_{\alpha} \sum_{i=1}^N \alpha_i - \frac{1}{2}\sum_{i=1}^N\sum_{j=1}^N \alpha_i\alpha_j y_i y_j x_i^T x_j \\
\text{subject to}~\sum_{i=1}^N \alpha_i y_i = 0,~\alpha_i \geq 0,~\forall i
\end{align}

\textbf{Decision Function:}
\begin{align}
f(x) = \sum_{i=1}^N \alpha_i y_i K(x_i, x) + w_0
\end{align}

\textbf{Common Kernels:}
\begin{itemize}[leftmargin=*,noitemsep,topsep=0pt]
    \item Linear: $K(x,z) = x^Tz$
    \item Polynomial: $K(x,z) = (1 + x^Tz)^q$
    \begin{itemize}[leftmargin=*,noitemsep]
        \item For $q=2$, includes all constant, linear, and quadratic terms
    \end{itemize}
    \item RBF: $K(x,z) = \exp(-\gamma\|x-z\|^2)$
    \begin{itemize}[leftmargin=*,noitemsep]
        \item Corresponds to infinite-dimensional feature space
        \item Parameter $\gamma$ controls width of Gaussian
    \end{itemize}
\end{itemize}

\textbf{Mercer's Condition:} A valid kernel function must produce a positive semi-definite Gram matrix.

\textbf{Kernel Construction Rules:}
\begin{itemize}[leftmargin=*,noitemsep,topsep=0pt]
    \item Sum: $K_1 + K_2$
    \item Product: $K_1 \cdot K_2$
    \item Scalar: $c \cdot K$
    \item Exponentiation: $\exp(K)$
\end{itemize}
\end{tcolorbox}

%%%%%%%%%%%%%%%%%%%%%%%%%%%%%%%%%%%%%%%%%%
\begin{tcolorbox}[breakable, title=6. Mixture Models Concepts, colback=mixture,boxsep=1pt,left=1pt,right=1pt,top=1pt,bottom=1pt]
\textbf{Mixture Models Overview:}
\begin{itemize}[leftmargin=*,noitemsep,topsep=0pt]
    \item Probabilistic models for data from multiple sources/distributions
    \item Each data point belongs to a latent (unobserved) class
    \item Models joint distribution p(x,z) of data x and latent class z
    \item Can be used for clustering, density estimation, and generative modeling
\end{itemize}

\textbf{General Form:}
\begin{align}
p(x) = \sum_{k=1}^K \pi_k p(x|z=k)
\end{align}
where $\pi_k$ are mixing coefficients ($\sum_k \pi_k = 1$) and $p(x|z=k)$ are component distributions.

\textbf{Estimation Challenge:}
\begin{itemize}[leftmargin=*,noitemsep,topsep=0pt]
    \item Maximum likelihood: $\log p(X) = \sum_{i=1}^N \log \sum_{k=1}^K \pi_k p(x_i|z=k)$
    \item Sum inside logarithm prevents direct optimization
    \item Solution: Use Expectation-Maximization (EM) algorithm
\end{itemize}

\textbf{Types of Mixture Models:}
\begin{itemize}[leftmargin=*,noitemsep,topsep=0pt]
    \item Gaussian Mixture Models (GMMs): Component distributions are Gaussian
    \item Mixture of Multinomials: For discrete data (text, genomics)
    \item Mixture of Experts: Different predictive models for different regions
    \item Latent Dirichlet Allocation: Documents as mixtures of topics
\end{itemize}
\end{tcolorbox}

%%%%%%%%%%%%%%%%%%%%%%%%%%%%%%%%%%%%%%%%%%
\begin{tcolorbox}[breakable, title=6.1 Expectation-Maximization (EM) Algorithm, colback=mixture,boxsep=1pt,left=1pt,right=1pt,top=1pt,bottom=1pt]
\textbf{When to use:} Use EM when you have latent variable models where direct maximum likelihood is intractable, especially for mixture models.

\textbf{General Procedure:}
\begin{enumerate}[leftmargin=*,noitemsep,topsep=0pt]
    \item Initialize parameters $\theta^{(0)}$
    \item Repeat until convergence:
    \begin{itemize}[leftmargin=*,noitemsep]
        \item E-step: Compute 
        $Q(\theta|\theta^{(t)}) = \mathbb{E}_{Z|X,\theta^{(t)}}[\log p(X, Z|\theta)]$
        \item M-step: $\theta^{(t+1)} = \arg\max_\theta Q(\theta|\theta^{(t)})$
    \end{itemize}
\end{enumerate}

\textbf{Theoretical Foundation:}
\begin{itemize}[leftmargin=*,noitemsep,topsep=0pt]
    \item EM maximizes a lower bound (ELBO) on the log likelihood:
    \begin{align}
    \log p(X|\theta) \geq \mathbb{E}_{Z|X,\theta^{(t)}}[\log p(X,Z|\theta)] - \mathbb{E}_{Z|X,\theta^{(t)}}[\log p(Z|X,\theta^{(t)})]
    \end{align}
    \item E-step: Compute posterior distribution of latent variables
    \item M-step: Maximize expected complete-data log likelihood
    \item Guaranteed to increase log likelihood at each iteration
\end{itemize}

\textbf{Mathematical Properties:}
\begin{itemize}[leftmargin=*,noitemsep,topsep=0pt]
    \item Converges to local maximum of likelihood
    \item ELBO becomes tight when using true posterior
    \item Different initializations may lead to different local maxima
\end{itemize}

\textbf{Applications Beyond Mixture Models:}
\begin{itemize}[leftmargin=*,noitemsep,topsep=0pt]
    \item Hidden Markov Models
    \item Latent factor models (PCA, factor analysis)
    \item Missing data imputation
\end{itemize}
\end{tcolorbox}

%%%%%%%%%%%%%%%%%%%%%%%%%%%%%%%%%%%%%%%%%%
\begin{tcolorbox}[breakable, title=6.2 Gaussian Mixture Models (GMM), colback=mixture,boxsep=1pt,left=1pt,right=1pt,top=1pt,bottom=1pt]
\textbf{When to use:} Use GMMs when your data appears to come from multiple Gaussian distributions or when you need soft clustering with uncertainty estimates.

\textbf{Model:} $p(x) = \sum_{k=1}^K \pi_k \mathcal{N}(x|\mu_k, \Sigma_k)$

\textbf{EM Algorithm for GMM:}
\begin{enumerate}[leftmargin=*,noitemsep,topsep=0pt]
    \item E-step: Calculate responsibilities (posterior probabilities)
    \begin{align}
    \gamma_{ik} = \frac{\pi_k \mathcal{N}(x_i|\mu_k, \Sigma_k)}{\sum_{j=1}^K \pi_j \mathcal{N}(x_i|\mu_j, \Sigma_j)}
    \end{align}
    
    \item M-step: Update parameters
    \begin{align}
    N_k &= \sum_{i=1}^N \gamma_{ik}\\
    \pi_k &= \frac{N_k}{N}\\
    \mu_k &= \frac{1}{N_k}\sum_{i=1}^N \gamma_{ik}x_i\\
    \Sigma_k &= \frac{1}{N_k}\sum_{i=1}^N \gamma_{ik}(x_i - \mu_k)(x_i - \mu_k)^T
    \end{align}
\end{enumerate}

\textbf{GMM vs. K-means:}
\begin{itemize}[leftmargin=*,noitemsep,topsep=0pt]
    \item GMM provides soft assignments (probabilities)
    \item GMM models covariance structure
    \item GMM is more flexible but requires more parameters
    \item K-means is a special case of GMM with:
    \begin{itemize}[leftmargin=*,noitemsep]
        \item Equal mixing coefficients $\pi_k = 1/K$
        \item Spherical covariances $\Sigma_k = \epsilon I$
        \item $\epsilon \to 0$ (hard assignments)
    \end{itemize}
\end{itemize}

\textbf{Choosing K:}
\begin{itemize}[leftmargin=*,noitemsep,topsep=0pt]
    \item BIC: $\text{BIC} = -2\log p(X|\hat{\theta}) + p\log N$
    \item AIC: $\text{AIC} = -2\log p(X|\hat{\theta}) + 2p$
    \item Cross-validation
\end{itemize}
\end{tcolorbox}

%%%%%%%%%%%%%%%%%%%%%%%%%%%%%%%%%%%%%%%%%%
\begin{tcolorbox}[breakable, title=6.3 Latent Dirichlet Allocation (LDA), colback=mixture,boxsep=1pt,left=1pt,right=1pt,top=1pt,bottom=1pt]
\textbf{When to use:} Use LDA for topic modeling in text data when you want to discover latent topics and their distributions across documents.

\textbf{Key Idea:} Documents are mixtures of topics, and topics are distributions over words.

\textbf{Generative Process:}
\begin{enumerate}[leftmargin=*,noitemsep,topsep=0pt]
    \item For each document:
    \begin{itemize}[leftmargin=*,noitemsep]
        \item Draw topic mixture $\theta_d \sim \text{Dir}(\alpha)$
    \end{itemize}
    \item For each topic:
    \begin{itemize}[leftmargin=*,noitemsep]
        \item Draw word distribution $\phi_k \sim \text{Dir}(\beta)$
    \end{itemize}
    \item For each word position in each document:
    \begin{itemize}[leftmargin=*,noitemsep]
        \item Draw topic $z_{d,n} \sim \text{Cat}(\theta_d)$
        \item Draw word $w_{d,n} \sim \text{Cat}(\phi_{z_{d,n}})$
    \end{itemize}
\end{enumerate}

\textbf{LDA as an Admixture Model:}
\begin{itemize}[leftmargin=*,noitemsep,topsep=0pt]
    \item Documents can belong to multiple topics
    \item Words within a document can come from different topics
    \item Extends mixture models with hierarchical structure
\end{itemize}

\textbf{Inference Methods:}
\begin{itemize}[leftmargin=*,noitemsep,topsep=0pt]
    \item Variational inference
    \item Gibbs sampling
    \item Both approximate posterior distributions for latent variables
\end{itemize}

\textbf{LDA Applications:}
\begin{itemize}[leftmargin=*,noitemsep,topsep=0pt]
    \item Document clustering
    \item Information retrieval
    \item Recommendation systems
    \item Feature extraction for text classification
\end{itemize}
\end{tcolorbox}

%%%%%%%%%%%%%%%%%%%%%%%%%%%%%%%%%%%%%%%%%%
\begin{tcolorbox}[breakable, title=9. Hidden Markov Models Concepts, colback=hmm,boxsep=1pt,left=1pt,right=1pt,top=1pt,bottom=1pt]
\textbf{HMM Overview:}
\begin{itemize}[leftmargin=*,noitemsep,topsep=0pt]
    \item Models sequential data with hidden states
    \item States follow Markov property (depends only on previous state)
    \item Each state generates an observation (emission)
    \item Useful for temporal or sequential data
\end{itemize}

\textbf{Key Components:}
\begin{itemize}[leftmargin=*,noitemsep,topsep=0pt]
    \item Hidden states $s_t$: Unobserved state at time $t$
    \item Emissions $x_t$: Observed values at time $t$
    \item Initial state distribution: $p(s_1)$
    \item Transition probabilities: $p(s_t|s_{t-1})$
    \item Emission probabilities: $p(x_t|s_t)$
\end{itemize}

\textbf{Common Problems:}
\begin{itemize}[leftmargin=*,noitemsep,topsep=0pt]
    \item Inference: Determine p(states|observations)
    \item Learning: Estimate parameters from observations
    \item Decoding: Find most likely state sequence
\end{itemize}

\textbf{Common Applications:}
\begin{itemize}[leftmargin=*,noitemsep,topsep=0pt]
    \item Speech recognition
    \item Natural language processing
    \item Biological sequence analysis
    \item Financial time series
\end{itemize}
\end{tcolorbox}

%%%%%%%%%%%%%%%%%%%%%%%%%%%%%%%%%%%%%%%%%%
\begin{tcolorbox}[breakable, title=9.1 Forward-Backward Algorithm for HMMs, colback=hmm,boxsep=1pt,left=1pt,right=1pt,top=1pt,bottom=1pt]
\textbf{When to use:} Use the Forward-Backward algorithm for efficient inference in HMMs, particularly for computing marginal probabilities of hidden states.

\textbf{Forward Pass ($\alpha$ values):}
\begin{align}
\alpha_t(s_t) &= p(x_1, ..., x_t, s_t)\\
\alpha_1(s_1) &= p(x_1|s_1)p(s_1)\\
\alpha_t(s_t) &= p(x_t|s_t) \sum_{s_{t-1}} p(s_t|s_{t-1})\alpha_{t-1}(s_{t-1})
\end{align}

\textbf{Backward Pass ($\beta$ values):}
\begin{align}
\beta_t(s_t) &= p(x_{t+1}, ..., x_n|s_t)\\
\beta_n(s_n) &= 1\\
\beta_t(s_t) &= \sum_{s_{t+1}} p(s_{t+1}|s_t)p(x_{t+1}|s_{t+1})\beta_{t+1}(s_{t+1})
\end{align}

\textbf{Inference with $\alpha$ and $\beta$:}
\begin{align}
p(x_1, ..., x_n) &= \sum_{s_t} \alpha_t(s_t)\beta_t(s_t)\\
p(s_t|x_1, ..., x_n) &\propto \alpha_t(s_t)\beta_t(s_t)\\
p(s_t, s_{t+1}|x_1, ..., x_n) &\propto \alpha_t(s_t)p(s_{t+1}|s_t)\\
&\quad\cdot p(x_{t+1}|s_{t+1})\beta_{t+1}(s_{t+1})
\end{align}

\textbf{Viterbi Algorithm (Best Path):}
\begin{enumerate}[leftmargin=*,noitemsep,topsep=0pt]
    \item Initialize:
    \begin{align}
    \delta_1(i) &= p(s_1 = i)p(x_1|s_1 = i)\\
    \psi_1(i) &= 0
    \end{align}
    
    \item Recursion:
    \begin{align}
    \delta_t(j) &= \max_i [\delta_{t-1}(i)p(s_t = j|s_{t-1} = i)]\\
    &\quad\cdot p(x_t|s_t = j)\\
    \psi_t(j) &= \arg\max_i [\delta_{t-1}(i)p(s_t = j|s_{t-1} = i)]
    \end{align}
    
    \item Termination:
    \begin{align}
    p^* &= \max_i [\delta_T(i)]\\
    s^*_T &= \arg\max_i [\delta_T(i)]
    \end{align}
    
    \item Backtracking:
    $s^*_t = \psi_{t+1}(s^*_{t+1})$ for $t = T-1,...,1$
\end{enumerate}

\textbf{HMM Training with EM:}
\begin{itemize}[leftmargin=*,noitemsep,topsep=0pt]
    \item E-step: Compute expected state occupancy and transitions using forward-backward
    \item M-step: Update parameters using these expectations
\end{itemize}
\end{tcolorbox}

%%%%%%%%%%%%%%%%%%%%%%%%%%%%%%%%%%%%%%%%%%
\begin{tcolorbox}[breakable, title=10. Markov Decision Processes Concepts, colback=mdp,boxsep=1pt,left=1pt,right=1pt,top=1pt,bottom=1pt]
\textbf{MDP Overview:}
\begin{itemize}[leftmargin=*,noitemsep,topsep=0pt]
    \item Framework for decision-making under uncertainty
    \item Agent takes actions in states to maximize rewards
    \item Environment transitions between states based on actions
    \item Goal: Find optimal policy (mapping from states to actions)
\end{itemize}

\textbf{Key Components:}
\begin{itemize}[leftmargin=*,noitemsep,topsep=0pt]
    \item States $S$: Possible situations for the agent
    \item Actions $A$: Possible choices in each state
    \item Transition function $p(s'|s,a)$: Probability of next state
    \item Reward function $r(s,a)$: Immediate reward for action
    \item Discount factor $\gamma$: Value of future rewards
    \item Policy $\pi(a|s)$: Strategy for choosing actions
\end{itemize}

\textbf{Value Functions:}
\begin{itemize}[leftmargin=*,noitemsep,topsep=0pt]
    \item State value $V^\pi(s)$: Expected return starting from state $s$, following policy $\pi$
    \item Action value $Q^\pi(s,a)$: Expected return starting from state $s$, taking action $a$, then following policy $\pi$
\end{itemize}

\textbf{Planning Horizons:}
\begin{itemize}[leftmargin=*,noitemsep,topsep=0pt]
    \item Finite horizon: Fixed number of time steps
    \item Infinite horizon: Unbounded time with discounting
\end{itemize}
\end{tcolorbox}

%%%%%%%%%%%%%%%%%%%%%%%%%%%%%%%%%%%%%%%%%%
\begin{tcolorbox}[breakable, title=10.1 MDP Planning Algorithms, colback=mdp,boxsep=1pt,left=1pt,right=1pt,top=1pt,bottom=1pt]
\textbf{When to use:} Use these planning algorithms when you have a fully defined MDP (states, actions, transitions, rewards) and need to find the optimal policy.

\textbf{Value Iteration:}
\begin{enumerate}[leftmargin=*,noitemsep,topsep=0pt]
    \item Initialize $V(s) = 0$ for all states
    \item Repeat until convergence:
    \begin{align}
    V'(s) &\leftarrow \max_a \left[r(s, a) + \gamma \sum_{s'} p(s'|s, a)V(s')\right]\\
    V(s) &\leftarrow V'(s)
    \end{align}
    \item Extract policy: 
    $\pi^*(s) = \arg\max_a \left[r(s, a) + \gamma \sum_{s'} p(s'|s, a)V^*(s')\right]$
\end{enumerate}

\textbf{Policy Iteration:}
\begin{enumerate}[leftmargin=*,noitemsep,topsep=0pt]
    \item Initialize policy $\pi$
    \item Repeat until convergence:
    \begin{itemize}[leftmargin=*,noitemsep]
        \item Policy Evaluation: Calculate $V^\pi(s)$
        \begin{align}
        V^\pi(s) = r(s, \pi(s)) + \gamma \sum_{s'} p(s'|s, \pi(s))V^\pi(s')
        \end{align}
        \item Policy Improvement:
        \begin{align}
        \pi'(s) &\leftarrow \arg\max_a \left[r(s, a) + \gamma \sum_{s'} p(s'|s, a)V^\pi(s')\right]\\
        \pi(s) &\leftarrow \pi'(s)
        \end{align}
    \end{itemize}
\end{enumerate}

\textbf{Value Iteration vs. Policy Iteration:}
\begin{itemize}[leftmargin=*,noitemsep,topsep=0pt]
    \item Value Iteration: More iterations, less computation per iteration
    \item Policy Iteration: Fewer iterations, more computation per iteration
    \item Value Iteration: O($|S|^2|A|$) per iteration
    \item Policy Evaluation: O($|S|^3$) for exact solution
\end{itemize}

\textbf{Bellman Equations:}
\begin{itemize}[leftmargin=*,noitemsep,topsep=0pt]
    \item Bellman Expectation Equation:
    $V^\pi(s) = \sum_a \pi(a|s)[r(s,a) + \gamma\sum_{s'} p(s'|s,a)V^\pi(s')]$
    \item Bellman Optimality Equation:
    $V^*(s) = \max_a [r(s,a) + \gamma\sum_{s'} p(s'|s,a)V^*(s')]$
\end{itemize}
\end{tcolorbox}

%%%%%%%%%%%%%%%%%%%%%%%%%%%%%%%%%%%%%%%%%%
\begin{tcolorbox}[breakable, title=Time Complexity Analysis, colback=white,boxsep=1pt,left=1pt,right=1pt,top=1pt,bottom=1pt]
\textbf{How to Calculate Time Complexity:}

\textbf{1. Identify basic operations}
\begin{itemize}[leftmargin=*,noitemsep,topsep=0pt]
    \item Arithmetic operations: O(1)
    \item Array/list access: O(1)
    \item Loop iterations: O(n) for n iterations
    \item Nested loops: O(n·m) for outer loop n times, inner loop m times
\end{itemize}

\textbf{2. Count operations in terms of input size}
\begin{itemize}[leftmargin=*,noitemsep,topsep=0pt]
    \item N = number of data points
    \item D = dimensionality of data
    \item K = number of clusters/classes
    \item T = number of iterations
\end{itemize}

\textbf{3. Focus on dominant terms}
\begin{itemize}[leftmargin=*,noitemsep,topsep=0pt]
    \item Ignore constants and lower-order terms
    \item Consider worst-case scenario
\end{itemize}

\textbf{Example Analysis: K-means}
\begin{enumerate}[leftmargin=*,noitemsep,topsep=0pt]
    \item For each iteration (T iterations):
    \begin{itemize}[leftmargin=*,noitemsep]
        \item For each data point (N points):
            \begin{itemize}[leftmargin=*,noitemsep]
                \item Calculate distance to each centroid (K centroids):
                    \begin{itemize}[leftmargin=*,noitemsep]
                        \item Each distance calculation: O(D) operations
                    \end{itemize}
            \end{itemize}
        \item Update K centroids: O(N·D) operations
    \end{itemize}
    \item Total complexity: O(T·N·K·D)
\end{enumerate}

\textbf{Common Algorithm Complexities:}
\begin{itemize}[leftmargin=*,noitemsep,topsep=0pt]
    \item Linear Regression: O(Nd² + d³) for optimal solution
    \item Gradient Descent: O(Ndt) for t iterations
    \item K-means: O(NKdt) for t iterations
    \item Decision Trees: O(dN log N) for training
    \item SVMs: O(N² to N³) depending on implementation
    \item Neural Networks: $O(N_t \cdot \sum_{\ell} n_{\ell} n_{\ell-1})$ for training
    \item PCA: $O(Nd^2 + d^3)$ for computation
    \item HMM Forward-Backward: $O(TK^2)$ for sequence length $T$, $K$ states
\end{itemize}
\end{tcolorbox}

\end{multicols}
\end{document}